\paragraph{Declare the company book before applying its methods.}
The preceding native examples use a fixed company context and chart of accounts.
A reusable calendar does not remove that restriction. A company may keep its cash
under account \texttt{101}, another under \texttt{bank.alpha}; their contract
records may also have different structures. The optional
\texttt{CompanyBook/1} reference makes those choices explicit before a method
reads balances or proposes an entry. It declares the entity, book, currency,
conditional case, period, account classifications and selected cash account.
Balances use the existing debit-positive convention and sum to zero. Each scalar
in the method's memory has a declared type and an address within a closed schema.

The payment example separates the accounting relation from those local names.
One synthetic company maps the cash and liability roles to
\texttt{bank.alpha} and \texttt{note.alpha}, with the outstanding amount under
\texttt{loan}. Another maps the roles to \texttt{101} and \texttt{230}, with the
amount under \texttt{contracts/obligation}. Both start with cash of USD~100 and an
outstanding liability of USD~30. A payment of USD~3.50 credits cash and debits
the liability by the same amount, leaving cash of USD~96.50 and a liability of
USD~26.50. The contract record falls to USD~26.50 and its cumulative payments
rise to USD~3.50. A second payment of USD~26.50 extinguishes the liability and
leaves cash of USD~70. Other accounts and unrelated contract records are unchanged.
The example assumes sufficient cash and a positive payment no larger than the
outstanding amount; the generic book does not impose a universal funding policy.

\texttt{CompanyBookMethod/1} declares an injective mapping from account roles to
the book's accounts, prerequisite predicates, monetary calculations, memory
replacements and postconditions. Bindings select declared account balances,
typed memory values or supplied event inputs. Their company, dimensions and
scale must agree exactly with the method's input declarations. A dollar amount
cannot become an inventory quantity merely because both are rational numbers.
The existing monetary interpreter derives the journal and separate gross cash
flows. The memory interpreter then replaces the specified scalar fields; this
profile rejects dynamic paths, new collection members and container writes.
Nested prerequisite and monetary definitions carry the enclosing method's
identity and version. A literal memory replacement is interpreted in the
destination's declared type.

\texttt{CompanyBookTransition/1} checks the entire request before assessing its
prerequisites. It derives monetary effects, stages memory changes and constructs
a new book from the same starting state. Only then does it evaluate the
postconditions against that proposed book. In the payment case, a postcondition
requires the credit balance of the liability account to equal the outstanding
contract amount. A failed postcondition returns no applied book, although the
staged calculation remains available for diagnosis. An unknown applicability
finding remains unresolved; a false finding can make the chosen method
unsupported. An inactive guarded prerequisite does not, by itself, cancel the
economic event: a separate decision determines whether an event should occur.

The request identifies both the starting book and the complete method by digest.
Application recomputes the proposal and compares its canonical JSON representation,
including Boolean and numeric types. Changing a saved journal, memory value or
authorization flag therefore cannot pass merely by replacing its digest. Event
identities cannot repeat, dates cannot precede the current book, and each event
must fall within the declared period. This is an atomic construction of a
provisional value; the separate posting-authorization and durable-storage rules
still govern any admitted or persisted transaction.

Five comparisons connect this reference to the retained native methods. Two
current-tax payments and three complete-lot sales reproduce their native journals,
gross flows and entire resulting states through explicit mappings. The mappings
retain every native field outside the declared changes; both payment cases also
reproduce the complete native financial reports. These checks concern the selected
synthetic cases. General reporting and closing, source admission, dynamic contract
collections, currency conversion and the broader corporate package still require
their own methods. The declared book supplies a common place to compose those
methods without weakening the older profiles' invariants.
